\documentclass{article}
\usepackage[T1]{fontenc}
\usepackage{lmodern}
\usepackage{graphicx}
\usepackage{amsmath,amssymb}
\usepackage[a4paper,margin=2cm]{geometry}
\usepackage[absolute,overlay]{textpos}
\setlength{\TPHorizModule}{1mm}
\setlength{\TPVertModule}{1mm}
\usepackage[indonesian]{babel}
\usepackage{hyperref}
\setlength{\emergencystretch}{2em}
% unit: Halaman 15

\begin{document}

% === Halaman 15 (llm) ===

\textbf{Contoh 3:} Misalkan barisan \textit{superincreasing} adalah $\{2, 3, 6, 13, 27, 52\}$, dan $m = 105$, dan $n = 31$.

Barisan \textit{non-superincreasing} (atau normal) \textit{knapsack} dihitung sbb:

\begin{align*}
2 \cdot 31 \pmod{105} &= 62 \\
3 \cdot 31 \pmod{105} &= 93 \\
6 \cdot 31 \pmod{105} &= 81 \\
13 \cdot 31 \pmod{105} &= 88 \\
27 \cdot 31 \pmod{105} &= 102 \\
52 \cdot 31 \pmod{105} &= 37
\end{align*}

Jadi, kunci publik adalah $\{62, 93, 81, 88, 102, 37\}$, sedangkan kunci privat adalah $\{2, 3, 6, 13, 27, 52\}$.

\end{document}
